%!TEX root = ../main.tex
% Resumen en español (~250 palabras). Las palabras clave se definen
% en main.tex con \palabrasclave{...}
%
% Guía de la cátedra (Entrega 0): el resumen responde, en este orden,
%   1. Problema: qué situación concreta de TI se aborda y a quién afecta.
%   2. Enfoque: qué se construye o investiga y con qué método o tecnología.
%   3. Valor: qué aporte original ofrece frente a lo existente.
% Impersonal, sin citas ni siglas sin expandir. Se reescribe al final del
% proyecto para que refleje lo que efectivamente se logró.

\begin{resumen}

Este proyecto nace como una red a lo largo del país, con una posibilidad
de expandirse, para conectar el aprendizaje y el conocimiento, elevando
la tradición del apoyo escolar en los hogares argentinos y brindándoles
nuevas herramientas. La iniciativa busca facilitarle a cualquier
estudiante, desde el nivel primario hasta el universitario, encontrar el
apoyo académico ideal, permitiendo que el saber circule sin barreras
geográficas o físicas.

Actualmente, el mercado de clases particulares opera de manera muy
fragmentada. Si bien la utilización de redes de vínculos entre personas
genera confianza, tiene un alcance muy limitado y se restringe a
entornos geográficos inmediatos. El proyecto centraliza la oferta en un
ecosistema digital donde la búsqueda, la reserva y los pagos conviven en
un solo sitio. Esto profesionaliza la labor docente y ofrece a las
familias un entorno transparente para potenciar los resultados
escolares.

Un pilar central es el diseño de un modelo de confianza para el
resguardo de los alumnos. Se plantea como objetivo deseado la
implementación de procesos de certificación que idealmente incluyan
validación de identidad, títulos y antecedentes de los tutores. Esta
meta busca garantizar que el vínculo educativo sea seguro y cumpla con
estándares de excelencia verificables.

Finalmente, el proyecto apunta a expandir el alcance de muchos
profesionales y a facilitar el acceso al apoyo escolar personalizado
mediante una plataforma digital con reputación dinámica y aulas
virtuales integradas. Al poner en valor la labor de los tutores y
formalizar el sector, este proyecto se posiciona como una
infraestructura esencial para conectar a los estudiantes con educadores
calificados, superando las barreras geográficas y generando un sustento
económico para estos últimos. El impacto de esta solución se observará
en la eficiencia del motor de matching, la estabilidad de las
conexiones remotas y el nivel de satisfacción de los usuarios.

\end{resumen}
